\documentclass[11pt,a4paper]{article}
\usepackage[T1]{fontenc}
\usepackage[utf8]{inputenc}
\usepackage{lmodern,amsmath,amssymb}
\usepackage[margin=25mm]{geometry}
\usepackage[hidelinks]{hyperref}
\hypersetup{pdftitle={Thermal retention and trigger discrimination},pdfauthor={navstokgap research working notes}}
\newcommand{\tightlist}{\setlength{\itemsep}{0pt}\setlength{\parskip}{0pt}}
\setlength{\emergencystretch}{3em}
\setcounter{secnumdepth}{0}
\begin{document}
\hypertarget{thermal-reliability-selects-a-detector-cost-not-a-universal-action}{%
\section{Thermal reliability selects a detector cost, not a universal
action}\label{thermal-reliability-selects-a-detector-cost-not-a-universal-action}}

A receiver with a fixed activated-rate prefactor cannot be made
arbitrarily ready while remaining quiet for a prescribed duration and
responding quickly to a signal. In the model below, those two
requirements impose an exact lower bound on signal-induced barrier
reduction. Interpreting that reduction as incident energy needs an
additional coupling premise. The bound retains bath temperature, timing
and error tolerances; it does not select a quantum action.

This is Q10's exploratory effective-model test, with a written
consistency review and no accepted-ledger promotion. The exponential
escape law and its signal coupling are physical hypotheses. They are not
a derivation of escape from Q09's conservative double well or a
microscopic model of photon detection.

\hypertarget{a-receiver-with-a-specified-retention-requirement}{%
\subsection{1. A receiver with a specified retention
requirement}\label{a-receiver-with-a-specified-retention-requirement}}

Let \(T>0\) be bath temperature, \(k_B\) Boltzmann's constant, and
\(\nu>0\) an attempt frequency in inverse seconds. A prepared metastable
receiver has residual activation barrier \(b\ge0\), measured in energy
units. Its first recorded escape is an absorbing event. The stipulated
dark rate is

\[r_0=\nu e^{-b/(k_BT)}. \tag{1}\]

Thus the survival probability S solves \(dS/dt=-r_0S\), \(S(0)=1\),
giving \(S(t)=\exp(-r_0t)\). Absorption describes a monitored first
event; maintaining a permanent physical record and resetting it remain
separate apparatus tasks. Retention here means absence of a false first
event while armed.

During a signal gate of duration \(t_g>0\) the barrier is reduced by an
amount \(u\ge0\), held constant throughout the gate. Use the explicitly
capped model

\[r_1=\nu e^{-\max(b-u,0)/(k_BT)}. \tag{2}\]

Here \(u\) denotes an energy change, whereas \(\nu\) denotes a rate. The
cap prevents an unphysical extrapolation to rates above the stipulated
attempt frequency. It is a model definition, not a Kramers formula near
a vanishing barrier. The same \(\nu\) in (1) and (2), an unchanged bath
and exponential holding times are essential premises. Sustaining u may
require work from a control source; no equality between u and input
pulse energy has yet been assumed.

Require false-event probability at most \(\epsilon\) over a dark
interval \(\tau>0\), and miss probability at most \(\eta\) in a
subsequent signal gate, conditional on surviving to that gate. For
\(0<\epsilon,\eta<1\), define

\[a=-\log(1-\epsilon)>0,\qquad c=\log(1/\eta)>0. \tag{3}\]

The two requirements are exactly

\[r_0\tau\le a,\qquad r_1t_g\ge c. \tag{4}\]

Conditional efficiency matters: the probability of both surviving the
dark interval and firing in the gate is
\(\exp(-r_0\tau)[1-\exp(-r_1t_g)]\). The separate bounds ensure it is at
least \((1-\epsilon)(1-\eta)\).

\hypertarget{exact-feasibility-and-the-optimal-residual-barrier}{%
\subsection{2. Exact feasibility and the optimal residual
barrier}\label{exact-feasibility-and-the-optimal-residual-barrier}}

Timely detection is impossible in this model if \(\nu t_g<c\), even with
zero remaining barrier. If \(\nu t_g\ge c\), define

\[B=\max\!\left(0,k_BT\log\frac{\nu\tau}{a}\right),\qquad
D=k_BT\log\frac{\nu t_g}{c}\ge0. \tag{5}\]

The feasible designs are precisely

\[b\ge B,\qquad u\ge\max(0,b-D). \tag{6}\]

Indeed, taking logarithms of (4) first gives
\(b\ge k_BT\log(\nu\tau/a)\). The signal inequality gives
\(\max(b-u,0)\le D\). Since \(D\ge0\), this is equivalent to
\(b-u\le D\), proving (6). Every point in (6) satisfies both original
probability bounds, including the equality cases. Minimizing over
\(b\ge0\) therefore yields

\[u_{\min}=\max(0,B-D)
 =k_BT\left[\log\frac{c\tau}{a t_g}\right]_+,
 \qquad \nu t_g\ge c, \tag{7}\]

where \([x]_+=\max(x,0)\). For the last equality, if \(\nu\tau\le a\)
then \(B=0\); feasibility also gives \(c\tau/(at_g)\le\nu\tau/a\le1\).
If \(\nu\tau>a\), subtract the two logarithms in (5). This checks both
branches without extending a negative barrier.

In the discriminating regime \(c\tau>at_g\), choose \(b=B\) and
\(u=B-D\). Both probability constraints then saturate. If
\(c\tau\le at_g\), a design with \(u=0\) can already satisfy them: the
timing/error specification is too weak to require signal discrimination.
A nonzero temperature alone is insufficient.

\hypertarget{what-becomes-of-preloading-and-a-slow-receiver}{%
\subsection{3. What becomes of preloading and a slow
receiver?}\label{what-becomes-of-preloading-and-a-slow-receiver}}

At fixed \(T,\nu,\tau,\epsilon\) with \(\nu\tau>a\), preloading cannot
reduce b below B while satisfying dark retention. This closes the
zero-residual-barrier route within the specified metastable class. A
receiver with bare barrier \(\Delta\) and a sustainable preload \(e_0\),
modeled by \(b=\Delta-e_0\), must satisfy

\[0\le e_0\le\Delta-B. \tag{8}\]

If \(\Delta<B\) the design cannot meet retention. A particle initially
at rest near Q09's unstable saddle is not automatically such a
sustainable metastable preload: it moves even without noise. Equation
(8) applies only if an actual preparation/control mechanism realizes (1)
for the full dark interval.

Reducing \(\nu\) alone improves dark survival but slows detection. The
gate floor \(\nu\ge c/t_g\) prevents using \(\nu\to0\) at fixed gate and
miss tolerance. In the positive-bound regime the same prefactor cancels
from (7), exposing the required contrast \(r_1/r_0\ge c\tau/(at_g)\).
Retention alone would not give this protection: it is the combined dark
and signal requirement that does so.

If the signal can instead change the prefactor, a rate-only alternative
is

\[r_0=\nu_0 e^{-b/(k_BT)},\qquad
r_1=\nu_1 e^{-b/(k_BT)}. \tag{9}\]

For any fixed \(b\ge0\) choose \(\nu_0=(a/\tau)\exp(b/(k_BT))\) and
\(\nu_1=(c/t_g)\exp(b/(k_BT))\). Both specifications hold with \(u=0\).
This is an exact counterexample to an inference from probabilities alone
to barrier reduction, not a construction of a zero-work autonomous
detector. Changing mobility, coupling or access to a reaction channel
has its own physical control cost, which (9) does not specify.

\hypertarget{energy-action-and-the-limits-of-the-bound}{%
\subsection{4. Energy, action and the limits of the
bound}\label{energy-action-and-the-limits-of-the-bound}}

To infer incident work W, add the explicit transduction premise

\[u\le\chi W,\qquad 0<\chi<\infty, \tag{10}\]

with fixed dimensionless \(\chi\) and with all sources contributing to
barrier control included in W. Then (7) implies

\[W\ge\frac{k_BT}{\chi}
 \left[\log\frac{c\tau}{a t_g}\right]_+. \tag{11}\]

A stored-energy amplifier need not obey (10) for signal work alone. Nor
is a barrier-height change automatically dissipated heat or reset work:
it concerns the energy landscape along an escape path, not the work
along the actual control trajectory. Q09's cycle accounting remains
necessary.

For Q08's narrow-band wave input of angular frequency \(\omega\),
canonical wave action is approximately \(W/\omega\). Under (10) its
corresponding lower bound is the right side of (11) divided by omega. A
finite-pulse wave-action definition must retain Q08's bandwidth
qualification. Alternatively \(Wt_g\) has action units but describes a
different quantity. Neither choice removes T, error rates, timing,
transduction or frequency from the result.

Within the stipulated rate-model class, \(T\to0\) at fixed \(\nu\) and
fixed specifications sends the optimal b, u and W bound to zero while
preserving all rate ratios. No microscopic low-temperature extrapolation
is claimed. Allowing larger \(\chi\) also removes a bound on signal
work; allowing \(t_g\ge c\tau/a\) removes the positive discrimination
requirement. These are explicit class limits, not quantum predictions.

\hypertarget{strategic-consequence-and-evidence-map}{%
\subsection{5. Strategic consequence and evidence
map}\label{strategic-consequence-and-evidence-map}}

Thermal reliability supplies a conditional cost for a ready detector
once a response deadline and a fixed signal coupling are included. It
does not by itself identify a universal action or yield exclusive
two-receiver events. Park further Arrhenius/threshold variants. The next
useful construction is an autonomous shared-resource receiver: can a
finite mechanical latch make two classical outputs compete for one
stored excitation while retaining reliable response? Specify its
resource, locality and output probabilities. That tests Q09's
exclusivity obstruction directly and can fail by requiring communication
or a prepared shared resource; it must not assume a Born rule.

Source context is limited. The existing
\href{../docs/batches/B23/harmonic-receiver-source-companion.md}{B23
companion}, Zwanzig p.~219 (21)--(24), supplies a bath-memory
construction, not (1). One discovery search located Hänggi, Talkner and
Borkovec, \emph{Reaction-rate theory: fifty years after Kramers},
Rev.~Mod. Phys. 62, 251--341 (1990),
\href{https://doi.org/10.1103/RevModPhys.62.251}{DOI}. Its bibliographic
metadata/abstract identify activated escape as established context. The
PDF was opened but no rate formula is imported from it; this is not a
passage or prior-art audit. Equations (1), (2) and (10) are declared
premises; (3)--(8) and (11) are written consequences; (9) tests an
omitted coupling premise. Independent proof and bounded literature
review remain required before ledger promotion. No computational
verification was used.
\end{document}
